\chapter{Discussion and Conclusions}\label{ch:conclusions}

\TODO{These are only some ideas that need a lot of work}

This book has developed a rigorous mathematical formalization of \gls{GDL}, establishing the theoretical foundations that unify modern deep learning architectures under a common geometric paradigm. Across four chapters, an arc has been constructed that starts from traditional neural networks, establishes the geometric framework of the ``5~Gs'', and shows how convolution and attention operations extend to non-Euclidean domains while preserving fundamental symmetries. The following sections synthesize the results obtained in relation to the four objectives set out in \autoref{sec:objetivos_trabajo}, identify the limitations that emerge from the theoretical development itself, and propose future directions grounded in the content of the book.


\section{Synthesis of Results and Contributions}

\subsection{Theoretical Foundations: The Framework of the 5~Gs}

The first objective (developing the mathematical framework of the ``5~Gs'') has been addressed in Chapters~\ref{ch:neural_networks} and~\ref{ch:gdl}. In \autoref{ch:neural_networks} the notion of an artificial neural network has been formalized starting from the definition of a neuron (\autoref{def:neurona_artificial}), subsequently establishing the foundations of measure theory and functional analysis needed to prove the universal approximation theorem (\autoref{sec:aproximacion_universal}). This result, which guarantees the expressive capacity of single-hidden-layer networks, constitutes the theoretical starting point from which the more complex architectures are justified. \autoref{ch:neural_networks} has shown the limitations inherent to traditional architectures: their inability to exploit data symmetries and the resulting parametric inefficiency.

\autoref{ch:gdl} has developed the complete framework of the 5~Gs (\autoref{subsec:cinco_gs}): differentiable manifolds as geometric domains, Lie groups as the algebraic description of symmetries, graphs as discretizations of these domains, geodesics as a mechanism for information propagation, and gauge theory to handle local invariances across coordinate systems. The key results of representation theory (the Peter-Weyl theorem (\autoref{thm:peter_weyl}) and Schur's lemma (\autoref{lem:schur})) provide the tools to decompose group representations into irreducible components, which is essential for constructing equivariant operators. The notions of equivariance (\autoref{def:equivariance_gdl}) and invariance (\autoref{def:invariance_gdl}) have been formalized as compatibility properties between functions and group actions, establishing the guiding principle of architectural design in \gls{GDL}. Gauge theory, developed through principal bundles (\autoref{def:fibrado_principal_formal}), has extended this principle to the case of local symmetries.

\subsection{Architectural Unification}

The second objective (showing that seemingly disparate architectures emerge as particular cases of the geometric paradigm) has been attained through the concept of equivariant \textit{message passing} as a unifying paradigm (\autoref{def:mpnn_unificador}). It has been formally shown that:

\begin{itemize}
	\item Traditional neural networks constitute \glspl{MPNN} over complete graphs (\autoref{prop:ann_mpnn}), where each neuron exchanges information with all the others.
	\item \glspl{CNN} perform \gls{message_passing} over regular grids (\autoref{thm:equivalencia_constructiva_cnn_gnn}), with shared weights determined by the translation group.
	\item \glspl{Transformer} implement adaptive \gls{message_passing} (\autoref{def:self_attention_basica}), where the communication weights are computed dynamically from the data.
\end{itemize}

This unification is not merely conceptual. The Weisfeiler-Leman test (\autoref{thm:wl_mpnn_equivalencia}) establishes precise expressive limits for \glspl{MPNN}, showing that their graph-distinguishing power coincides exactly with that of the combinatorial 1-WL algorithm. The generalization hierarchy (\autoref{rem:jerarquia_generalizacion}) formalizes the inclusion relations between the various architectural families, providing a complete theoretical map of their relative capabilities.

\subsection{Geometric Extensions}

The third objective (analyzing the extensions of convolutions and attention mechanisms to non-Euclidean domains) has been developed in Chapters~\ref{ch:cnn} and~\ref{ch:attention}.

\autoref{ch:cnn} has constructed a progression of convolutional generalizations of increasing complexity: group convolutions, steerable networks, spherical convolutions, $\text{E}(3)$/$SE(3)$-equivariant networks, gauge equivariant networks, and Clifford-valued convolutions. This chain culminates in the equivariant universality theorem (\autoref{theorem:equivariant_universality_cnns}), which guarantees that equivariant \glspl{CNN} can approximate any continuous equivariant function, thus extending the classical universal approximation result to the geometric context. The topological extensions, developed through the Hodge Laplacians (\autoref{def:hodge_laplacians_cnn}), have provided tools for operating over higher-order structures (simplicial and cellular complexes).

\autoref{ch:attention} has extended attention mechanisms to diverse geometric domains. Equivariant attention via Clifford algebra (\autoref{sec:clifford_attention}) allows operating directly on multivectors, preserving the symmetries of geometric algebra. Hyperbolic attention (\autoref{def:atencion_hiperbolica}) extends attention mechanisms to spaces of constant negative curvature, suitable for data with hierarchical structure. Harmonic analysis has acted as a unifying tool throughout both chapters, connecting the spatial and spectral perspectives of geometric operations.

\subsection{Theory-Architecture Connection}

The fourth objective (illustrating how the formalized geometric principles manifest in modern models) has been addressed through the analysis of representative architectures. AlphaFold~\citep{jumper2021alphafold} exemplifies $SE(3)$ equivariance in protein structure prediction. DeepSets and PointNet illustrate permutation invariance, with the corresponding universality theorem (\autoref{thm:deepsets_universal}) guaranteeing their expressive capacity over set functions. \gls{GATr} demonstrate the applicability of Clifford algebras in attention architectures.

The \gls{CNN}-attention equivalence table (\autoref{tab:cnn_attention_equivalencias_fundamentales}) synthesizes the formal correspondences between both paradigms, showing that they share mathematical structure despite their apparent differences. The correspondence theorem (\autoref{thm:correspondencia_cnn_attention}) and the convergence conjecture (\autoref{conj:convergencia_cnn_attention}) formalize this relation, suggesting that under appropriate conditions the convolutional and attention operations converge to the same operator.


\section{Limitations and Open Problems}

\subsection{Expressive Power and the Weisfeiler-Leman Barrier}

The equivalence between \glspl{MPNN} and the 1-WL test (\autoref{thm:wl_mpnn_equivalencia}) establishes a fundamental expressive barrier: standard \glspl{MPNN} cannot distinguish graphs that the 1-WL test does not distinguish. The $k$-WL extensions overcome this limitation, but at a computational cost that grows exponentially with $k$, which makes them impractical for graphs of moderate size. A central open problem is the design of architectures with polynomial cost that surpass the 1-WL barrier, possibly through the use of additional structural features or higher-order attention mechanisms.

The barrier itself has been refined in recent years. For discrete graphs, \citet{zhang2023biconnectivity} show that no architecture of the 1-WL class can decide biconnectivity (cut vertices and cut edges), a property that a message passing scheme over generalized distances recovers at polynomial cost, and \citet{zhang2024beyondwl} replace the binary comparison with the $k$-WL hierarchy by \textit{homomorphism expressivity}, which characterizes exactly which substructure counts each architectural family can compute; for geometric graphs, the geometric WL test of \citet{joshi2023expressive} (\autoref{subsubsec:wl_geometrico}) shows that invariant layers are strictly less expressive than equivariant ones and that the expressive power depends on the depth, the tensor order, and the body order of the messages. Over-squashing, the complementary obstruction, is now understood quantitatively: \citet{digiovanni2023oversquashing} bound the sensitivity of a node representation to the input features of a distant node by a quantity governed by the effective resistance (equivalently, the commute time) between the two nodes, so that the topology of the graph, rather than the depth or width of the network, determines which pairs of nodes can exchange information.

\subsection{Tension between Rigidity and Adaptability}

The trade-off between rigidity and adaptability (\autoref{rem:tradeoff_rigidez_adaptabilidad}), formalized in \autoref{ch:attention}, identifies a fundamental tension in the design of geometric architectures. Equivariant \glspl{CNN} incorporate symmetries in a rigid manner, which provides data efficiency and theoretical guarantees, but limits their adaptability when the symmetries are only approximate. Attention mechanisms offer flexibility to learn context-dependent relations, but at the cost of greater data demand and weaker structural regularization. The \gls{CNN}-attention convergence conjecture (\autoref{conj:convergencia_cnn_attention}) suggests that both paradigms converge under appropriate conditions, but the optimal point on the rigidity-adaptability spectrum for each concrete domain remains an open problem.

A third option, which decouples equivariance from the architecture, has emerged between the two poles. Frame averaging \citep{puny2022frameaveraging} makes an arbitrary backbone $f$ exactly equivariant by averaging it over a small input-dependent set of group elements $F(\mathbf{x}) \subset G$ (a \textit{frame}, satisfying $F(g\mathbf{x}) = g\,F(\mathbf{x})$), $\Phi(\mathbf{x}) = \frac{1}{|F(\mathbf{x})|}\sum_{g \in F(\mathbf{x})} \rho_{\mathcal{Y}}(g)\, f(g^{-1}\mathbf{x})$, without constraining the layers of $f$ and without losing universality; learned canonicalization \citep{kaba2023canonicalization} reduces the frame to a single element predicted by an auxiliary equivariant network that brings the input to a canonical pose (see \autoref{subsec:canonicalizacion_frames}). When the group itself is unknown, symmetry discovery methods such as LieGAN \citep{yang2023liegan} learn the generators of a Lie algebra of symmetries from the data by adversarial training, supplying the missing ingredient for the geometric meta-learner of \autoref{def:meta_learner_geometrico}.

\subsection{Gauge Equivariance and Global Topology}

The gauge theory developed in \autoref{ch:cnn} through principal bundles (\autoref{def:fibrado_principal_formal}) provides a rigorous framework for local symmetries. However, the global topological obstructions (nontrivial bundles, nontrivial holonomy around cycles) pose both theoretical and computational difficulties. The gap between the elegance of the gauge formalism and the implementation of efficient architectures at scale constitutes a significant practical challenge: current gauge equivariant \glspl{CNN} require careful discretizations that may not preserve the global topological properties of the underlying space.

\subsection{Higher-Order Structures}

The topological tools developed in this book (Hodge Laplacians (\autoref{def:hodge_laplacians_cnn}), persistent homology, simplicial complexes) provide a mathematically rigorous framework for capturing higher-order relations. Nevertheless, the architectures that operate over these structures are at an early stage of maturity compared to architectures over graphs. The design of efficient convolution and attention operations over general cellular complexes, which preserve the algebraic properties of the underlying topology without incurring prohibitive computational costs, remains a relevant open problem. The position paper by \citet{papamarkou2024tdlposition}, signed by a large part of the community, lists the open problems of \gls{TDL} in these terms: an expressivity theory over complexes analogous to the WL hierarchy for graphs, the scalability of message passing over high-dimensional cells, the lack of benchmarks where the topological structure is demonstrably necessary, and the connection with cellular sheaves and other algebraic structures; the survey by \citet{papillon2023tdlarchitectures} organizes the existing architectures under a single message passing template over hypergraphs and simplicial, cellular, and combinatorial complexes, which makes a systematic comparison between them possible.


\section{Future Directions}

The following lines of research are grounded directly in the results and limitations identified in this book.

\subsection{Hybrid CNN-Attention Architectures}

The correspondence theorem (\autoref{thm:correspondencia_cnn_attention}) and the convergence conjecture (\autoref{conj:convergencia_cnn_attention}) between convolutions and attention, together with the analysis of the rigidity-adaptability spectrum (\autoref{rem:tradeoff_rigidez_adaptabilidad}), suggest that a natural direction is the design of architectures that combine both paradigms in a controlled manner. The theoretical framework established in \autoref{ch:attention} provides the formal tools to interpolate between the rigidity of equivariant convolutions and the adaptability of attention, potentially combining the data efficiency of the former with the expressiveness of the latter.

Two concrete architectural patterns illustrate this interpolation. The first combines a convolutional kernel, which processes local features with geometric guarantees, with attention heads that capture adaptive global dependencies, mediating between the two through a learnable switching mechanism.

\begin{definition}[ConvNext-Attention Hybrid]\label{def:convnext_attention}
A hybrid architecture that combines:
\begin{itemize}
    \item \textbf{Convolutional kernel}: To process local features with geometric guarantees
    \item \textbf{Attention heads}: To capture adaptive global dependencies
    \item \textbf{Dynamic switching}: A learnable mechanism that decides when to use each component
\end{itemize}
\end{definition}

The second carries the retentive mechanism over to geometric domains, replacing the positional decay of sequences with a decay over the geodesic distance of the domain:

\begin{definition}[Geometric Retentive Networks]\label{def:retentive_networks_geometricas}
Extension of the retentive mechanism to geometric domains:
\begin{equation}
    \text{Retention}(Q, K, V) = Q \odot \sum_{j} \text{decay}(d_\mathcal{M}(i,j)) \cdot K_j \odot V_j
\end{equation}
where $d_\mathcal{M}$ is the geodesic distance in the geometric domain $\mathcal{M}$.
\end{definition}

\subsection{Complex Geometric Domains}

The framework of the 5~Gs (\autoref{subsec:cinco_gs}) admits natural extensions to product spaces that combine geometries of different curvature (hyperbolic, spherical, and Euclidean) in a single model. These structures arise naturally in data with simultaneous hierarchical and cyclic components. In the topological direction, the results on simplicial complexes and cellular sheaves of \autoref{ch:cnn} provide a basis for the development of \textit{sheaf neural networks}, which would generalize graph networks by incorporating the algebraic structure of sheaves over cellular complexes.

A third extension concerns the symmetries themselves rather than the domain. The framework of the 5~Gs assumes a global symmetry group, whereas most real domains offer only partial symmetries: local isomorphisms between neighborhoods of a graph, or radial isometries between patches of a manifold. Groupoids (\autoref{subsubsec:grupoides}) provide the algebraic structure for such symmetries, and recent work establishes equivariant convolutions and universal approximation theorems over Lie groupoids \citep{astwood2026liegroupoid, maruyama2025categorical}, while the analysis of fibration symmetries \citep{velarde2026fibration} shows that the local, non-automorphic symmetries of a graph already constrain which functions a message passing network can compute. Whether these formalisms yield architectures that outperform Natural Graph Networks \citep{dehaan2020natural} and gauge equivariant networks in practice remains open: the theory is in place, but no benchmarks exist yet.

\subsection{Generative Models on Manifolds and Lie Groups}\label{subsec:modelos_generativos_variedades}

The manifold and Lie group material of \autoref{ch:gdl} (the Riemannian metric (\autoref{def:metrica_riemanniana}), geodesics (\autoref{def:geodesica}), the exponential map (\autoref{def:mapa_exponencial_riemanniano}), and the Laplace-Beltrami operator (\autoref{def:laplace_beltrami})) has been used in this book to construct discriminative equivariant operators. The same tools are exactly what is needed to model probability distributions on non-Euclidean spaces, a direction that the book has not covered and that has become one of the most active applications of \gls{GDL}. The only missing ingredient is stochastic: Brownian motion on a manifold, the diffusion whose generator is $\tfrac{1}{2}\Delta_\mathcal{M}$. Three constructions illustrate the direction.

\textbf{Riemannian score-based generative modelling.} \citet{debortoli2022riemanniansgm} extend score-based diffusion models to a compact Riemannian manifold $(\mathcal{M}, g)$ by replacing the Euclidean noising process with Riemannian Brownian motion $\mathrm{d}\mathbf{X}_t = \mathrm{d}\mathbf{B}^\mathcal{M}_t$, whose transition density $p_t(\cdot \mid \mathbf{x}_0)$ is the heat kernel of the Laplace-Beltrami operator, $\partial_t p_t = \tfrac{1}{2}\Delta_\mathcal{M} p_t$, and which converges to the uniform distribution with respect to the Riemannian volume. A network $s_\theta(t, \mathbf{x}) \in T_\mathbf{x}\mathcal{M}$ is trained by score matching to approximate the Riemannian score $\nabla_\mathbf{x} \log p_t(\mathbf{x})$, a tangent vector field, and samples are generated by integrating the time-reversed SDE $\mathrm{d}\mathbf{Y}_t = s_\theta(T - t, \mathbf{Y}_t)\,\mathrm{d}t + \mathrm{d}\mathbf{B}^\mathcal{M}_t$, discretized as a geodesic random walk $\mathbf{Y}_{k+1} = \exp_{\mathbf{Y}_k}\!\big(\gamma\, s_\theta(T - t_k, \mathbf{Y}_k) + \sqrt{\gamma}\,\boldsymbol{\xi}_k\big)$ with $\boldsymbol{\xi}_k$ Gaussian in $T_{\mathbf{Y}_k}\mathcal{M}$. On symmetric spaces such as the sphere the heat kernel is an explicit series in the eigenfunctions of $\Delta_\mathcal{M}$ (the spherical harmonics), which makes the training target computable; on general manifolds it must be approximated, and this is the main practical cost of the method.

\textbf{Riemannian flow matching.} \citet{chen2024flowmatchinggeometries} avoid stochastic simulation altogether. Given a source distribution $p_0$ on $\mathcal{M}$ (typically uniform) and the data distribution $p_1$, one draws $\mathbf{x}_0 \sim p_0$, $\mathbf{x}_1 \sim p_1$ and connects them by the geodesic interpolant
\begin{equation}
    \mathbf{x}_t = \exp_{\mathbf{x}_0}\!\big(t \log_{\mathbf{x}_0}(\mathbf{x}_1)\big), \qquad t \in [0, 1],
\end{equation}
whose velocity $\dot{\mathbf{x}}_t = \log_{\mathbf{x}_t}(\mathbf{x}_1)/(1 - t)$ is the tangent vector pointing from $\mathbf{x}_t$ toward the data point. A neural vector field $v_\theta(t, \cdot)$ tangent to $\mathcal{M}$ is trained by regression onto this velocity,
\begin{equation}\label{eq:riemannian_flow_matching}
    \mathcal{L}_{\text{RFM}}(\theta) = \mathbb{E}_{t \sim \mathcal{U}[0,1],\; \mathbf{x}_0 \sim p_0,\; \mathbf{x}_1 \sim p_1}\, \big\| v_\theta(t, \mathbf{x}_t) - \dot{\mathbf{x}}_t \big\|^2_{g(\mathbf{x}_t)},
\end{equation}
where $\|\cdot\|_{g}$ is the norm induced by the Riemannian metric. Although the regression target is conditional on the pair $(\mathbf{x}_0, \mathbf{x}_1)$, the minimizer is the marginal vector field that transports $p_0$ to $p_1$, and sampling reduces to integrating the ODE $\dot{\mathbf{x}} = v_\theta(t, \mathbf{x})$ on the manifold with exponential map steps. Training is simulation-free whenever $\exp$ and $\log$ are available in closed form: the sphere (\autoref{ex:mapa_exp_esfera}), the flat torus, hyperbolic space, and $SO(3)$ via Rodrigues' formula (\autoref{ex:mapa_exp_so3}). On general geometries (meshes with boundaries, for instance), the geodesic distance in the interpolant is replaced by a premetric computed from a spectral decomposition of the Laplace-Beltrami operator, at the price of approximating the target vector field.

\textbf{$SE(3)$ diffusion for protein backbones.} A backbone of $N$ residues is represented, following the AlphaFold convention, by $N$ rigid frames $\mathbf{T}_i = (\mathbf{R}_i, \mathbf{x}_i) \in SE(3)$, where the rotation orients the local triad of backbone atoms and the translation places the $\text{C}_\alpha$ atom; a backbone is thus a point of the Lie group $SE(3)^N$. \citet{yim2023se3diffusion} build the diffusion directly from the group structure $SE(3) = SO(3) \ltimes \mathbb{R}^3$: with a suitable left-invariant metric, Brownian motion on $SE(3)$ decouples into a Euclidean Brownian motion on the translations (after removing the centre of mass, so that a stationary distribution exists) and a Brownian motion on $SO(3)$ with its bi-invariant metric, whose transition density is the isotropic Gaussian on $SO(3)$.

\begin{definition}[Isotropic Gaussian on $SO(3)$]\label{def:igso3}
The isotropic Gaussian distribution $\mathcal{IG}_{SO(3)}(\mathbf{R}; \mathbf{R}_0, \sigma^2)$ is the transition density, with respect to the normalized Haar measure, of Brownian motion on $SO(3)$ with the bi-invariant metric, started at $\mathbf{R}_0$ and run for time $\sigma^2$. It depends only on the rotation angle $\omega \in [0, \pi]$ of the relative rotation $\mathbf{R}_0^\top \mathbf{R}$, given by $\cos\omega = (\operatorname{tr}(\mathbf{R}_0^\top \mathbf{R}) - 1)/2$, through
\begin{equation}\label{eq:igso3_densidad}
    f(\omega; \sigma^2) = \sum_{\ell = 0}^{\infty} (2\ell + 1)\, e^{-\ell(\ell+1)\sigma^2/2}\, \frac{\sin\big((\ell + \tfrac{1}{2})\omega\big)}{\sin(\omega/2)}.
\end{equation}
\end{definition}

The series in \eqref{eq:igso3_densidad} is the Peter-Weyl expansion (\autoref{thm:peter_weyl}) of a class function on $SO(3)$: $\chi_\ell(\omega) = \sin((\ell + \tfrac{1}{2})\omega)/\sin(\omega/2)$ is the character of the irreducible representation of dimension $2\ell + 1$, and $-\ell(\ell+1)$ is the eigenvalue of the Laplace-Beltrami operator on the corresponding isotypic component. The heat kernel of any compact Lie group has the same form, $\sum_\rho d_\rho\, e^{-\lambda_\rho t}\, \chi_\rho$, which is why Brownian motion on such a group is computable from its character table alone. The score of \eqref{eq:igso3_densidad} is obtained by differentiating the series in $\omega$, the score network outputs one element of the Lie algebra $\mathfrak{so}(3)$ per residue, mapped back to the group by the exponential map, and the network is made $SE(3)$-equivariant so that the generated distribution over backbones is $SE(3)$-invariant: an equivariant vector field started from an invariant distribution transports it to an invariant one.

This last point is where the open problem lies. Building equivariance into the generative model, as above, yields exact invariance of the sample distribution and data efficiency; AlphaFold~3 \citep{abramson2024alphafold3} instead discarded the $SE(3)$-equivariant structure module of its predecessor in favour of a diffusion module that operates on raw atomic coordinates without frames or equivariant processing, and recovers the symmetry only approximately through random rotations and translations of the training data. Whether, and at what scale of data and model, augmentation makes constrained equivariant generative models unnecessary is the generative counterpart of the rigidity-adaptability tension discussed in \autoref{ch:attention} (\autoref{rem:tradeoff_rigidez_adaptabilidad}) and of the equivariance-scalability dilemma analysed in \autoref{subsubsec:dilema_equivarianza_escalabilidad}.

\subsection{Generalized Harmonic Analysis}

The central role of harmonic analysis in this book (from spectral convolutions to spherical harmonics and the Chebyshev decomposition of \autoref{ch:cnn}) suggests the exploration of extensions to broader classes of groups. The Peter-Weyl theorem (\autoref{thm:peter_weyl}), the cornerstone of group convolutions, requires compactness of the group; extending the framework to non-compact groups (such as the Lorentz group or the affine groups) demands tools of non-commutative harmonic analysis that go beyond those developed in this book. An advance in this direction would broaden the scope of applicability of equivariant \glspl{CNN} to symmetries relevant in relativistic physics and general signal processing.


\subsection{Toward a Unified Theory}\label{subsec:hacia_teoria_unificada}

Two recent syntheses map the terrain that such a theory would have to reconcile: \citet{saezdeocarizborde2025mathfoundations} organize \gls{GDL} from the side of the model, around the mathematical structures (symmetry groups, manifolds, graphs) imposed on the architecture, whereas \citet{papillon2024beyondeuclid} organize it from the side of the data, classifying methods by the geometric, topological, and algebraic structure of their inputs; a unified theory should explain why the two taxonomies coincide where they do. Beyond the design of concrete architectures, the \gls{CNN}-attention convergence conjecture (\autoref{conj:convergencia_cnn_attention}) points toward a unified theory whose contours are beginning to emerge. A first line concerns the limits of expressivity: if both paradigms converge to the same limit set of operators, one may expect their approximation cost to converge as well.

\begin{conjecture}[Universal Limits of Geometric Approximation]\label{conj:limites_universales}
For every continuous function $f: \mathcal{M} \to \mathbb{R}$ over a compact manifold $\mathcal{M}$, and every $\epsilon > 0$, there exist CNN and Attention architectures of complexities $C_{CNN}(\epsilon)$ and $C_{Att}(\epsilon)$ such that:
\begin{equation}
    \frac{C_{Att}(\epsilon)}{C_{CNN}(\epsilon)} \to 1 \quad \text{as } \epsilon \to 0
\end{equation}
This conjecture suggests that the difference in architectural complexity vanishes in the limit of perfect approximation.
\end{conjecture}

A second line seeks a properly geometric information theory, one that weights statistical uncertainty by the geometry of the underlying domain.

\begin{definition}[Geometric Fisher Entropy]\label{def:entropia_geometrica_fisher}
For a distribution $p$ over a Riemannian manifold $(\mathcal{M}, g)$, the geometric entropy is defined as:
\begin{equation}
    H_g(p) = -\int_\mathcal{M} p(x) \log p(x) \sqrt{\det g(x)} \, dx
\end{equation}
This quantity captures both the statistical uncertainty and the underlying geometric complexity.
\end{definition}

\begin{conjecture}[Principle of Minimal Geometric Information]\label{conj:principio_informacion_minima}
The optimal GDL architectures minimize the geometric information necessary to represent the target function:
\begin{equation}
    \mathcal{A}^* = \arg\min_\mathcal{A} \left[ \mathcal{L}(\mathcal{A}) + \lambda H_g(\text{representations}_\mathcal{A}) \right]
\end{equation}
where $\mathcal{L}$ is the empirical loss and $\lambda$ controls the information-precision trade-off.
\end{conjecture}

A third line, complementary to the previous ones, is geometric meta-learning: rather than fixing the symmetry group a priori, the system itself discovers it from the data and builds the corresponding equivariant architecture.

\begin{definition}[Geometric Meta-Learner]\label{def:meta_learner_geometrico}
A system that automatically learns the relevant symmetries of a domain:
\begin{equation}
    \text{Meta-GDL}: \mathcal{D} \mapsto G_{\text{optimal}} \mapsto \mathcal{A}_{G_{\text{optimal}}}
\end{equation}
where $\mathcal{D}$ is the data, $G_{\text{optimal}}$ is the discovered group of symmetries, and $\mathcal{A}_{G_{\text{optimal}}}$ is the resulting architecture.
\end{definition}

These lines converge on a principle of synthesis: for each geometric problem, with a known structure and a required degree of adaptability, there would exist an optimal hybrid architecture that balances both ends of the rigidity-adaptability spectrum.

\begin{remark}[Principle of Architectural Synthesis]\label{rem:principio_sintesis}
For every geometric problem with known structure $\mathcal{S}$ and required adaptability $\mathcal{A}$, there exists an optimal hybrid architecture:
\begin{equation}
    \mathcal{H}^* = \arg\min_{\mathcal{H}} \left[ \mathcal{L}(\mathcal{H}) + \lambda_1 d(\mathcal{H}, \mathcal{S}) + \lambda_2 (1 - \text{Adaptability}(\mathcal{H})) \right]
\end{equation}
where $d(\mathcal{H}, \mathcal{S})$ measures the distance to the known structure and the parameters $\lambda_i$ balance the constraints.
\end{remark}


\section{Final Reflection}

The development presented in this book shows that geometry is not merely a tool for improving the performance of deep learning models, but an organizing principle that allows explaining why certain architectures work for certain types of data. From this deep theoretical perspective, the distinction between the ``convolutional'' and the ``attentional'' becomes increasingly less relevant: both paradigms are distinct points on a single continuum of geometric operators, and the future of the field lies less in their competition than in their principled synthesis. The synthesis of \autoref{ch:gdl} (\autoref{sec:sintesis_marco_unificador}) has shown that the differences between \glspl{ANN}, \glspl{CNN}, \glspl{GNN}, and \glspl{Transformer} reflect choices about three fundamental aspects: the geometric domain, the symmetry group, and the connectivity structure.

The mathematical framework that underpins this unification (Lie theory, Riemannian geometry, gauge theory, algebraic topology) possesses a power that transcends the mere retrospective description of existing architectures. The equivariant universality results and the convergence theorems between paradigms provide tools with predictive capacity: they allow anticipating which geometric constraints are necessary for a given domain and which expressive guarantees the resulting architecture will offer. In this sense, \gls{GDL} constitutes not only a taxonomy of current architectures, but a theoretical framework for the principled design of future ones.
